% Section 11.2, translated from sv/chapters/11/exercises.tex.

\section{Exercises}
\label{c11:sec:uppgifter}

\begin{aside}
\textbf{How to work with the exercises.} The exercises review Chapters~\ref{c7:ch}--\ref{c10:ch}
before the practice test and are best done on your own, one part for each chapter. Start with the
parts you feel least sure of.

\facitnotis{L11}
\end{aside}

% ------------------------------------------------------------------------------------------
\uppgiftsdel{Part 1}{Vectors}
These exercises review \kapref{7}.

\uppgift{1.1}{Calculating with vectors}
The vectors $\mathbf{u} = (2, 5)$ and $\mathbf{v} = (-3, 1)$ are given. Find:
\begin{delar}(4)
  \task $\mathbf{u} + \mathbf{v}$
  \task $\mathbf{u} - \mathbf{v}$
  \task $3\mathbf{u}$
  \task $2\mathbf{u} - 3\mathbf{v}$
\end{delar}

\uppgift{1.2}{Magnitude and angle}
\begin{parts}
  \item Calculate $\abs{(8, 6)}$.
  \item Calculate $\abs{(-5, -12)}$.
  \item Find the angle to the positive $x$-axis for $(8, 6)$.
  \item Find the angle for $(-5, -12)$. Remember the quadrant correction.
\end{parts}

\uppgift{1.3}{Scalar product}
\begin{parts}
  \item Calculate $\mathbf{u} \cdot \mathbf{v}$ for $\mathbf{u} = (3, 1)$ and
    $\mathbf{v} = (2, -6)$.
  \item What does the result say about the vectors?
  \item Calculate the angle between $\mathbf{u} = (1, 2)$ and $\mathbf{v} = (3, 1)$.
\end{parts}

\uppgift{1.4}{Unit vector and a given length}
The vector $\mathbf{v} = (-8, 6)$ is given.
\begin{parts}
  \item Find the unit vector in the direction of $\mathbf{v}$.
  \item Find a vector of length $25$ in the same direction as $\mathbf{v}$.
  \item Find a vector of length $5$ in the opposite direction to $\mathbf{v}$.
\end{parts}

\uppgift{1.5}{Voltages as phasors}
The voltages across a resistor and a capacitor in series are
$\mathbf{U}_R = (9, 0)\,\text{V}$ and $\mathbf{U}_C = (0, -12)\,\text{V}$.
\begin{parts}
  \item Calculate the total voltage $\mathbf{U} = \mathbf{U}_R + \mathbf{U}_C$.
  \item Calculate $\abs{\mathbf{U}}$.
  \item Calculate the angle to the positive $x$-axis.
  \item Why is $\abs{\mathbf{U}}$ not equal to $9 + 12 = 21\,\text{V}$?
\end{parts}

% ------------------------------------------------------------------------------------------
\uppgiftsdel{Part 2}{Functions}
These exercises review \kapref{8}.

\uppgift{2.1}{Domain and range}
Find the domain $D_f$ and the range $V_f$ of:
\begin{delar}(3)
  \task $f(x) = \sqrt{x - 2}$
  \task $f(x) = \dfrac{1}{x + 1}$
  \task $f(x) = -x^2 + 4$
\end{delar}

\uppgift{2.2}{A linear model of a thermistor}
The resistance of a thermistor depends linearly on the temperature as $R(T) = kT + m$.
Measurements give $R(0^{\circ}\text{C}) = 200\,\Omega$ and $R(40^{\circ}\text{C}) = 120\,\Omega$.
\begin{parts}
  \item Find $k$ and $m$.
  \item Calculate $R$ at $25^{\circ}\text{C}$.
  \item At what temperature is $R = 100\,\Omega$?
  \item Is the function increasing or decreasing?
\end{parts}

\uppgift{2.3}{An exponential model}
The number of faults in a system decreases by $15\,\%$ per month. At the start there are $80$
faults.
\begin{parts}
  \item Write a function $f(x)$ for the number of faults after $x$ months.
  \item Calculate the number of faults after $6$ months.
  \item After how many months are there fewer than $20$ faults?
\end{parts}

\uppgift{2.4}{Power as a power function}
The power in a resistor is given by $P(I) = 25I^2$ watts for $0 \leq I \leq 2\,\text{A}$.
\begin{parts}
  \item Calculate $P(0.4)$.
  \item Calculate $P(2)$.
  \item Find $V_P$ for the given interval.
  \item At what current is $P = 36\,\text{W}$?
\end{parts}

\uppgift{2.5}{A rational function}
The function $f(x) = \dfrac{x^2 - 9}{x - 3}$ is given.
\begin{parts}
  \item Factorise the numerator.
  \item Which values of $x$ are not permitted?
  \item Simplify the expression.
  \item Find the zero of the function.
\end{parts}

% ------------------------------------------------------------------------------------------
\uppgiftsdel{Part 3}{Trigonometry}
These exercises review \kapref{9}.

\uppgift{3.1}{Angle measure}
\begin{delar}(2)
  \task Convert $135^{\circ}$ to radians.
  \task Convert $-120^{\circ}$ to radians.
  \task Convert $\dfrac{5\pi}{6}$ to degrees.
  \task Convert $2\,\text{rad}$ to degrees.
\end{delar}

\uppgift{3.2}{Properties from the equation}
An AC voltage is given by $u(t) = 6\sin\left(400\pi t - \dfrac{\pi}{3}\right)$ volts.
\begin{parts}
  \item Give the amplitude.
  \item Give the angular frequency and the frequency.
  \item Give the period.
  \item How many milliseconds does the phase correspond to, and is the curve early or late?
\end{parts}

\uppgift{3.3}{Writing the equation}
An AC voltage has the amplitude $12\,\text{V}$, the period $T = 2\,\text{ms}$ and the phase
$30^{\circ}$. Find $u(t)$ in the form $\abs{U}\sin(\omega t + \delta)$ with the phase in radians.

\uppgift{3.4}{Instantaneous value, RMS value and power}
An AC voltage is given by $u(t) = 8\sin(100\pi t)$ volts.
\begin{parts}
  \item Calculate $u(5\,\text{ms})$.
  \item Calculate $U_{\text{RMS}}$.
  \item The voltage is applied across a resistor $R = 50\,\Omega$. Calculate the power.
\end{parts}

\uppgift{3.5}{A trigonometric equation}
An AC voltage is given by $u(t) = 10\sin(200\pi t)$ volts.
\begin{parts}
  \item At what time is $u(t) = 5\,\text{V}$ for the first time?
  \item Give the next time at which $u(t) = 5\,\text{V}$.
  \item At what time is the peak value first reached?
\end{parts}

% ------------------------------------------------------------------------------------------
\uppgiftsdel{Part 4}{Logarithms and decibels}
These exercises review \kapref{10}.

\uppgift{4.1}{Exponential and logarithmic equations}
Solve the equations.
\begin{delar}(4)
  \task $2^x = 128$
  \task $10^x = 500$
  \task $e^{2x} = 5$
  \task $\log(x - 1) = 2$
\end{delar}

\uppgift{4.2}{The laws of logarithms}
Calculate or simplify.
\begin{delar}(4)
  \task $\log 4 + \log 25$
  \task $\log 500 - \log 5$
  \task $3\log 2$
  \task $\ln e^3$
\end{delar}

\uppgift{4.3}{Gain in dB}
An amplifier has $U_{\text{in}} = 20\,\text{mV}$ and $U_{\text{out}} = 2\,\text{V}$.
\begin{parts}
  \item Calculate the gain in dB.
  \item After the amplifier comes a filter that attenuates $12\,\text{dB}$. Calculate the total
    gain in~dB.
  \item Calculate the total linear gain.
  \item Calculate the output voltage after the filter.
\end{parts}

\uppgift{4.4}{Power in dB and dBm}
The power level in dBm is given relative to $1\,\text{mW}$:
\[
  P_{\text{dBm}} = 10\log_{10}\left(\frac{P}{1\,\text{mW}}\right)
\]
\begin{parts}
  \item An amplifier has $P_{\text{in}} = 2\,\text{W}$ and $P_{\text{out}} = 50\,\text{W}$.
    Calculate the gain in dB.
  \item A transmitter has the output power $20\,\text{dBm}$. Give the power in mW.
  \item How many dBm does $1\,\text{W}$ correspond to?
\end{parts}

\uppgift{4.5}{Discharging a capacitor}
The voltage is given by $u(t) = 12e^{-t/(RC)}$ volts with $R = 22\,\text{k}\Omega$ and
$C = 100\,\mu\text{F}$.
\begin{parts}
  \item Calculate the time constant $\tau = RC$.
  \item Calculate the voltage after $3$ seconds.
  \item After how long has the voltage halved?
  \item After how long is the voltage down to $1\,\text{V}$?
\end{parts}
